% !TEX root = Online-main.tex

% \chuong[]{ÔN TẬP TỔNG HỢP}

\cau Cho $\triangle ABC$, lấy $M$ là trung điểm của $BC$, trên tia $AM$ lấy điểm $D$ sao cho $AM = MD$. Chứng minh tứ giác $ABCD$ là hình bình hành.
\dauchamK{12}

\cau Cho hình bình hành $ABCD$, hai đường chéo cắt nhau tại $O$. Lấy $M$, $N$ lần lượt là trung điểm của $OD$, $OB$. $E$ là giao điểm của $AM$ và $CD$, $F$ là giao điểm của $CN$ và $AB$.
\begin{tuluan}
    \item Chứng minh $AMCN$ là hình bình hành.
    \item Chứng minh $DE = BF$.
\end{tuluan}
\dauchamK{13}


\cau Cho hình bình hành $ABCD$. Gọi $E$, $K$ lần lượt là trung điểm của $CD$ và $AB$. Đường chéo $BD$ cắt $AE$, $AC$, $CK$ lần lượt tại $N$, $O$, $M$.
\begin{tuluan}
    \item Chứng minh $AECK$ là hình bình hành.
    \item Chứng minh 3 điểm $E, O, K$ thẳng hàng.
    \item Chứng minh $DN = NM = MB$.
    \item Chứng minh $AE = 3KM$.
\end{tuluan}
\dauchamE{1}

\cau Cho $\triangle ABC$ nhọn có $AB < AC$. Các đường cao $BE, CF$ cắt nhau tại $H$. Gọi $M$ là trung điểm của $BC$. Từ $B$ kẻ đường thẳng vuông góc với $AB$ và từ $C$ kẻ đường thẳng vuông góc với $AC$ hai đường thẳng này cắt nhau tại $K$.
\begin{tuluan}
    \item Chứng minh $BHCK$ là hình bình hành.
    \item Chứng minh $H, M, K$ thẳng hàng.
    \item Từ $H$ vẽ $HG \perp BC$. Trên tia $HG$ lấy $I$ sao cho $HG = GI$. Chứng minh tứ giác $BIKC$ là hình thang cân.
\end{tuluan}
\dauchamE{1}

\cau Cho $\triangle ABC$ vuông tại $A$ có $AB < AC$, đường cao $AH$ và trung tuyến $AE$. Gọi $D, F$ lần lượt là hình chiếu của $E$ trên $AB, AC$.
\begin{tuluan}
    \item Chứng minh $BDFE$ là hình bình hành.
    \item Chứng minh $DFEH$ là hình thang cân.
    \item Lấy $M$ sao cho $F$ là trung điểm của $EM$ và $N$ sao cho $F$ là trung điểm của $BN$. Chứng minh $A, N, M$ thẳng hàng.
\end{tuluan}
\dauchamE{1}

\cau Cho $\triangle ABC$ vuông tại $A$ có $AB < AC$, $M$ là trung điểm của $BC$. Kẻ $ME \perp AB$ ($E \in AB$), kẻ $MF \perp AC$ ($F \in AC$).
\begin{tuluan}
    \item Chứng minh $EF = \dfrac{BC}{2}$.
    \item Gọi $AK$ là đường cao của $\triangle ABC$. Chứng minh $KMFE$ là hình thang cân.
\end{tuluan}
\dauchamE{1}

\cau Cho $\triangle ABC$ vuông tại $A$ có $AB < AC$, đường cao $AH$. Từ $H$ kẻ $HM \perp AB$ ($M \in AB$), kẻ $HN \perp AC$ ($N \in AC$). Gọi $I$ là trung điểm của $HC$, lấy $K$ trên tia $AI$ sao cho $I$ là trung điểm của $AK$.

\begin{tuluan}
    \item Chứng minh $AC \parallel HK$.
    \item Chứng minh $MNCK$ là hình thang cân.
    \item $MN$ cắt $AH$ tại $O$, $CO$ cắt $AK$ tại $D$. Chứng minh $AK = 3AD$.
\end{tuluan}
\dauchamE{1}

\cau Cho $\triangle ABD$ vuông tại $A$ có $AB < AD$, $M$ là trung điểm của $BD$. Lấy $C$ sao cho $M$ là trung điểm của $AC$.
\begin{tuluan}
    \item Chứng minh $ABCD$ là hình chữ nhật.
    \item Trên tia đối của tia $DA$ lấy điểm $E$ sao cho $DA = DE$. Gọi $I$ là trung điểm của $CD$. Chứng minh $IB = IE$.
    \item Kẻ $AH \perp BD$. Lấy $K$ sao cho $H$ là trung điểm của $AK$. Chứng minh $BDCK$ là hình thang cân.
\end{tuluan}
\dauchamE{1}